\subsection{舵机观测通路异步化的可行性分析}\label{sec:async-read}

\subsubsection{上游 TODO：基于 txPacket/rxPacket 拆分的流水线方案}\label{sec:async-read-todo}

与异步写完全没有任何上游代码痕迹相对，
\texttt{huggingface/lerobot} 主仓库的
\texttt{src/lerobot/motors/motors\_bus.py}
在 \texttt{\_setup\_sync\_reader} 与 \texttt{sync\_write} 之间留有一段
针对``读''的 TODO 注释占位：

\begin{verbatim}
# TODO(aliberts, pkooij): Implementing something like this could
# get even much faster read times if need be.
# Would have to handle the logic of checking if a packet has been
# sent previously though but doable.
# This could be at the cost of increase latency between the moment
# the data is produced by the motors and the moment it is used by
# a policy.
# def _async_read(self, motor_ids, address, length):
#     if self.sync_reader.start_address != address or ...:
#         self._setup_sync_reader(motor_ids, address, length)
#     else:
#         self.sync_reader.rxPacket()
#         self.sync_reader.txPacket()
#     for id_ in motor_ids:
#         value = self.sync_reader.getData(id_, address, length)
\end{verbatim}

注释明确指向一种``流水线（pipelining）''方案：把当前 \texttt{txRxPacket()}
中的``发请求''与``收回包''两个原本绑在一起的步骤拆开，
让每次 \texttt{\_async\_read()} 调用先\textbf{接收上一次请求的回包}，
再\textbf{发起这一次的新请求}，下次调用再来收这一次的回包。
本质上是把单次串口往返的等待时间隐藏到两次相邻调用之间。

\subsubsection{与相机后台线程式异步的本质区别}\label{sec:async-read-not-thread}

需要特别澄清：上游规划的``异步读''与相机的 \texttt{async\_read}
\textbf{完全不是一个机制}。
\begin{itemize}[leftmargin=2em]
  \item 相机 \texttt{async\_read}：由独立后台线程持续 \texttt{cv2.VideoCapture.read()}，
        主线程仅取最新帧。属于\textbf{多线程并发}模式，
        受益于每台相机独占 USB 端口，没有总线竞争问题。
  \item 舵机 \texttt{\_async\_read}（TODO）：\textbf{完全在主线程内执行}，
        没有任何后台线程。仅仅是把单次原子的 \texttt{txRxPacket()}
        拆成两次连续调用之间的``收 + 发''重叠。属于\textbf{串行流水线}模式，
        因此不会与 \texttt{sync\_write} 在总线上撞车。
\end{itemize}

也正因为是流水线而非并发，4.2 节中关于半双工总线竞争（约束 A）和
错误暴露延迟（约束 C）的论证对舵机异步读\textbf{并不适用}：
异步读不引入并发，不需要锁，错误依然在主线程当帧抛出。

\subsubsection{落地代价：驱动改造与观测滞后}\label{sec:async-read-cost}

尽管异步读在原理上可行，要真正落地仍有两项代价：

\textbf{代价 1——需要修改 scservo\_sdk 驱动层}。
当前 Feetech SDK 把 \texttt{txPacket()} 与 \texttt{rxPacket()}
的调用封装在 \texttt{txRxPacket()} 里作为单个原子操作，
对外只暴露这一接口。要落地上游 TODO 的 pipelining 方案，
必须深入 scservo\_sdk 内部，把内部的 tx/rx 状态机暴露给 \texttt{motors\_bus.py}，
还需新增``上次发送的请求参数''缓存以及``参数不匹配则退化为同步读''的兼容路径。
这一改造涉及一个属于第三方厂商的 C/Python 混合驱动，
维护与上游升级跟踪都会有额外负担。

\textbf{代价 2——观测延迟 1 帧}。
TODO 注释明确写了：``This could be at the cost of increase latency between
the moment the data is produced by the motors and the moment it is used by a policy.''
异步读在第 $N$ 次调用拿到的实际上是第 $N-1$ 次请求时的电机状态。
对 ACT 这类闭环 policy 而言，
``此刻的观测''比物理时刻晚了一整帧（约 33 ms @30 FPS），
理论上会让控制误差略增。

\subsubsection{收益估算}\label{sec:async-read-roi}

按第三章的实测数据，单次 \texttt{sync\_read} 的串口往返耗时约 1.2 ms。
异步读最多能把这 1.2 ms 从主线程的关键路径上拿掉。然而：
\begin{itemize}[leftmargin=2em]
  \item 在当前 ACT 主导的循环中（推理约 2000 ms），
        节省 1.2 ms 对总循环时间的影响小于 0.1\%；
  \item 即便后续推理被加速到 100 ms 量级（INT8 + ONNX），
        异步读省下的 1.2 ms 也仅占约 1\%，仍属边际收益；
  \item 真正可能让异步读发挥作用的场景，是在没有视觉、
        没有大模型推理的纯实时控制系统里（控制频率 500 Hz 以上），
        而这类场景与本论文的 LeRobot ACT 用例并不重合。
\end{itemize}

\subsubsection{小结}\label{sec:async-read-conclusion}

舵机异步读在原理上可行（pipelining 不破坏半双工约束），
但综合来看 ROI 不足以支持改造工作量：
当前推理瓶颈下省下 1.2 ms 几乎不可见，
且需要改第三方 SDK 并接受 1 帧观测滞后。
本论文将其归入``结构上可行、当前不值得做''的优化项，
作为推理加速完成后可以重新评估的潜在方向，但不在本论文的实验范围内。
